\documentclass[a4paper,11pt]{article}
\usepackage[utf8]{inputenc}
\usepackage[T1]{fontenc}
\usepackage[french]{babel}
\usepackage{amsmath,amsfonts,amssymb}
\usepackage{geometry}
\usepackage{fancyhdr}
\usepackage{titlesec}
\usepackage{enumitem}
\usepackage{listings}

\geometry{margin=2.5cm}

% Header and footer configuration
\pagestyle{fancy}
\fancyhf{}
\fancyhead[L]{Licence 1 Informatique}
\fancyhead[R]{CERI - Avignon Université}
\fancyfoot[C]{Mathématiques Discrètes \hfill TD4 - Preuves de correction \hfill Page \thepage/4}

% Title formatting
\titleformat{\section}
{\Large\bfseries}
{\thesection}
{1em}
{}

\titleformat{\subsection}
{\large\bfseries}
{\thesubsection}
{1em}
{}

% Configuration listings pour pseudo-code
\lstset{
    basicstyle=\ttfamily\small,
    keywordstyle=\bfseries,
    commentstyle=\itshape\color{gray},
    frame=single,
    breaklines=true,
    showstringspaces=false
}

\begin{document}

\begin{center}
{\Huge\bfseries TD4 - Preuves de correction} \\[0.5cm]
{\large Hiver 2026} \\[0.3cm]
\end{center}

\vspace{0.5cm}

L'objectif de ce TD est de pratiquer les preuves de correction d'algorithmes.

\section{Procédure simple}

\begin{lstlisting}
y <- 2
z <- x + y
x <- y - z
\end{lstlisting}

\begin{enumerate}
\item Montrez la correction partielle de l'extrait ci-dessus pour l'assertion initiale $x = 5$ et l'assertion finale $x = 14$

\item Montrez la correction partielle de l'extrait ci-dessus pour l'assertion initiale $x = -5$ et l'assertion finale $x = -6$

\item Montrez la correction partielle de l'extrait ci-dessus pour l'assertion initiale $x = -2$ et l'assertion finale $x = 0$
\end{enumerate}

\section{Procédures avec blocs conditionnels}

\begin{enumerate}
\item Montrez que l'extrait de procédure suivant est partiellement correct pour l'assertion initiale $\top$ et l'assertion finale $z > 0$

\begin{lstlisting}
Si x < 0 Alors
    z <- 0
FinSi
\end{lstlisting}

\item Montrez que l'extrait de procédure suivant est partiellement correct pour l'assertion initiale $\top$ et l'assertion finale $(x \leq y \wedge min = x) \vee (x > y \wedge min = y)$

\begin{lstlisting}
Si x < y Alors
    min <- x
Sinon
    min <- y
FinSi
\end{lstlisting}
\end{enumerate}

\section{Procédures avec boucles}

\subsection{Calcul de puissance}

On souhaite montrer que la procédure suivante est correcte pour l'assertion initiale ``$x \in \mathbb{R}, n \in \mathbb{N}^{**}$'' et l'assertion finale ``\texttt{puissance}$(x,n) = x^n$''

\begin{lstlisting}
Procedure puissance(x : reel, n : entier strictement positif)
    puissance <- 1
    i <- 0
    TantQue i < n Faire
        puissance <- puissance * x
        i <- i + 1
    FinTantQue
    Renvoyer puissance
FinProcedure
\end{lstlisting}

\begin{enumerate}
\item Donnez un invariant de boucle

\item Vérifiez la correction partielle de la procédure en utilisant l'invariant de boucle de la question précédente

\item Montrez que la procédure termine et déduisez-en qu'elle est correcte
\end{enumerate}

\subsection{Division entière}

On souhaite montrer que la procédure suivante est correcte pour l'assertion initiale ``$a, d \in \mathbb{N}^*$'' et l'assertion finale ``$q, r \in \mathbb{N}^*$ où $a = dq + r$ et $0 \leq r < d$''

\begin{lstlisting}
Procedure division(a, d : entiers strictement positifs)
    q <- 0
    r <- a
    TantQue r >= d Faire
        r <- r - d
        q <- q + 1
    FinTantQue
    Renvoyer q, r
FinProcedure
\end{lstlisting}

\begin{enumerate}
\item Donnez un invariant de boucle

\item Vérifiez la correction partielle de la procédure en utilisant l'invariant de boucle de la question précédente

\item Montrez que la procédure termine et déduisez-en qu'elle est correcte
\end{enumerate}

\section{Procédures récursives}

\begin{enumerate}
\item Montrez que la procédure \texttt{somme}$(n)$ donnée calcule $\sum_{k=0}^{n} k$

\begin{lstlisting}
Procedure somme(n : entier positif)
    Si n = 0 Alors
        Renvoyer 0
    Sinon
        Renvoyer n + somme(n - 1)
    FinSi
FinProcedure
\end{lstlisting}

\item Montrez que la procédure \texttt{factorielle}$(n)$ donnée calcule $n!$

\begin{lstlisting}
Procedure factorielle(n : entier positif)
    Si n = 0 Alors
        Renvoyer 1
    Sinon
        Renvoyer n * factorielle(n - 1)
    FinSi
FinProcedure
\end{lstlisting}
\end{enumerate}

\section{Analyser des procédures}

\subsection{Procédure $p_{5,1}$}

\begin{lstlisting}
Procedure p_5,1(n : entier strictement positif, m : entier)
    Si n = 1 Alors
        Renvoyer m
    Sinon
        Renvoyer m + p_5,1(n - 1, m)
    FinSi
FinProcedure
\end{lstlisting}

\begin{enumerate}
\item Quel est le résultat renvoyé par la procédure $p_{5,1}$ en donnant en entrée $n = 2$ et $m = 2$

\item Quel est le résultat renvoyé par la procédure $p_{5,1}$ en donnant en entrée $n = 3$ et $m = 0$

\item Quel est le résultat renvoyé par la procédure $p_{5,1}$ en donnant en entrée $n = 3$ et $m = 4$

\item Déduisez des questions précédentes ce que calcule la procédure $p_{5,1}$

\item Montrez que la procédure $p_{5,1}$ est correcte
\end{enumerate}

\subsection{Procédure $p_{5,2}$}

\begin{lstlisting}
Procedure p_5,2(x : reel)
    y <- 2
    z <- x + y
    Si z > 0 Alors
        z <- 3
    Sinon
        z <- 0
    FinSi
    Renvoyer z
FinProcedure
\end{lstlisting}

\begin{enumerate}
\item Quel est le résultat renvoyé par la procédure $p_{5,2}$ en donnant en entrée $x = 3$

\item Quel est le résultat renvoyé par la procédure $p_{5,2}$ en donnant en entrée $x = -2$

\item Quel est le résultat renvoyé par la procédure $p_{5,2}$ en donnant en entrée $x = -4$

\item Déduisez des questions précédentes ce que calcule la procédure $p_{5,2}$

\item Montrez que la procédure $p_{5,2}$ est correcte
\end{enumerate}

\section{Identifier des règles d'inférences}

\subsection{Pour blocs conditionnels}

\begin{enumerate}
\item En utilisant la seconde règle d'inférence pour la correction partielle des blocs conditionnels vue en cours, proposez une règle d'inférence pour prouver la correction partielle du bloc conditionnel suivant :

\begin{lstlisting}
Si c_1 Alors
    S_1
SinonSi c_2 Alors
    S_2
Sinon
    S_3
FinSi
\end{lstlisting}

\item Généralisez pour un bloc conditionnel de la forme :

\begin{lstlisting}
Si c_1 Alors
    S_1
SinonSi c_2 Alors
    S_2
SinonSi c_3 Alors
    S_3
Sinon
    S_n
FinSi
\end{lstlisting}
\end{enumerate}

\subsection{Pour des cas quelconques}

\begin{enumerate}
\item On suppose que les hypothèses suivantes sont vraies :
\begin{itemize}
\item $q_0 \Rightarrow q_1$
\item $p\{S\}q_1$
\end{itemize}

Montrez qu'on peut en déduire $p\{S\}q$

\item On suppose que les hypothèses suivantes sont vraies :
\begin{itemize}
\item $p_0 \Rightarrow p_1$
\item $p_1\{S\}q$
\end{itemize}

Montrez qu'on peut en déduire $p_0\{S\}q$
\end{enumerate}

\end{document}
